\documentclass[conference]{IEEEtran}
\input preamble.tex
\input macros.tex
\input notation.tex
\graphicspath{{figs/}}   % use PDF/PNG  under ./figs
\linenumbers


\begin{document}



\title{Printing in a JIFFY: a parallel-in time heat transfer solver for additive manufacturing
%A parallel-in-time transient heat transfer for  additive manufacturing
  \vspace*{-30pt}
}

\maketitle
\begin{abstract}
Quality certification of mission-critical additively manufactured (AM) components can take more than a year due to the intricate coupling between the process parameters and the resulting part quality.  
This lengthy certification requires finding the right process parameters  and creates major bottlenecks in R\&D  across the aerospace, automotive, and biomedical device sectors.  
Predictive modeling and simulation offer a way to shorten this time frame. 
However, they are challenging because 3D printing spans very large ranges of spatial and temporal scales.

In this work, we isolate a single key physics challenge in 3D printing: predicting the temperature field.  
We introduce \method,  a highly specialized parallel-in-time solver for heat transfer simulations with a moving point source.
Given the source path and an approximate Green’s function, we construct a directed acyclic graph that encodes the space-time dependencies between sources and targets. We then decompose this graph into spatial path segments.  
These segments can be processed in parallel, and any remaining dependencies are handled through a small number of additional iterations that are local in both space and time.  
As a result, \method achieves nearly ideal weak scaling and very good strong scaling.  
We implement the solver in Python, with support for GPUs and distributed-memory parallelism via the Message Passing Interface (MPI).  
For instance, we can simulate a complex multilayer print on eight GPUs in under 30 minutes.  
The effective uniform-grid spatial resolution exceeds  100 billion points with 100s of thousands of time steps. 
A state-of-the-art single-GPU adaptively refined implementation that does not use temporal parallelism would require more than a day to perform the same computation.
\end{abstract}


\section{Introduction}
Many processes in metal additive manufacturing involve a moving heat source. 
Examples include direct energy deposition, welding, and laser power bed fusion (LPBF). 
To make ideas concrete, we describe \method in the context of LPBF.

Laser powder bed fusion (LPBF) is governed by a highly localized moving heat source that repeatedly melts and resolidifies the material along a prescribed scan path \cite{debroy2018additive, bikas2016additive}. The resulting thermal history determines melt-pool geometry, cooling behavior, and solidification conditions, which in turn influence defects, residual stresses, and final microstructure \cite{qin2023grainnn, qin2024graingnn, chadwick2021development, basak2016epitaxy}. For many process-design and microstructure-oriented studies, a transient thermal model captures the dominant quantities of interest, including the melt-pool shape, thermal gradient $G$, and solidification rate $R$, without requiring full melt-pool fluid dynamics or powder-scale multiphysics \cite{debroy2018additive,bikas2016additive}. 
LPBF is modeled as the standard nonlinear transient heat-transfer model for LPBF. 
Although the laser-matter interactions are highly complex, this formulation is widely used in metal AM because it can accurately predict quantities of interest like melt pool shape, $G$, and $R$~\cite{leonor2024go, keller2017application,debroy2018additive}. 

\paragraph{Problem statement}
Let $u(\mathbf{x},t)$ denote the non-dimensional temperature in the computational domain $\Omega \subset \mathbb{R}^3$. In this work, we consider the following heat equation
\begin{equation}
\frac{\partial u}{\partial t} = \Delta u + f_{\text{laser}},
\qquad \text{in } \Omega,
\label{eq:heat_nondim_intro}
\end{equation}
where $f_{\text{laser}}$ represents the moving laser heat source. 
Typically $\Omega$ is regular box as in \cref{fig:hermes_moving_domain}. 
On the top surface where the heating takes place, we have a Neumann boundary condition that involves the laser source (typically a Gaussian), as well as convective and radiative losses.
The remaining boundaries have zero fluxes (adiabatic). Finally, the initial condition is a uniform ambient temperature. 
The actual LPBF formulation involves nonlinearities in both specific heat and thermal conductivity coefficients (not depicted here). 
However, surprisingly, \method can be extended to the nonlinear case without major modifications because the nonlinearity is localized in the melt pool. See \cref{seq:conclusions} for  additional discussion. So, we restrict our attention to  \cref{eq:heat_nondim_intro}.

% Let $T(\mathbf{x},t)$ denote temperature in the computational domain $\Omega\subset\mathbb{R}^3$. The governing equation is
% \begin{equation}
% \rho c_{p}(T)\frac{\partial T}{\partial t}=\Div( K \Grad T) + f_{\text{laser}} \qquad \text{in } \Omega,
% \label{eq:full_heat_dim}
% \end{equation}
% where $\rho$ is the density, $K$ is the thermal conductivity, and $c_{p}(T)$ is an effective temperature-dependent heat capacity used to incorporate latent heat over the melting interval. On the top surface, heating is driven by a moving Gaussian laser source, $f_\text{laser}$, together with convective and radiative losses, while the remaining boundaries are taken to be adiabatic. The initial condition is a uniform ambient temperature. 

\begin{figure}[h]
\centering
\includegraphics[width=\linewidth]{HERMES_Demonstration2.png}
\caption{
Illustration of a representative state-of-the-art moving-domain SOTA solver such as GO-MELT or HERMES.  
Each box corresponds to the solution at a single time step.  
The computational box travels with the laser at every time step and must be sufficiently large to maintain accuracy.  
The solutions are temporally coupled because the initial condition for the box at time $t_{i}$ is obtained from the solution at time $t_{j}$.  
This chain of dependencies is typically irregular and hard to predict, which is why most SOTA solvers proceed sequentially in time.  
The inset in the upper right displays the temperature field and the associated solid–liquid interface.  
The region outlined in red is fully molten, and its subsequent solidification defines the final material geometry.  
Accurately predicting the temperature and especially its gradients within the melt pool is essential for accessing the quality of the 3D print. 
}
\label{fig:hermes_moving_domain}
\end{figure}


\paragraph{Multiscale challenges}
Despite its seemingly simple form,~\cref{eq:heat_nondim_intro} remains difficult to solve because of the vast spatiotemporal scales involved.
Melt pools typically measure on the order of tens to hundreds of microns, while the manufactured components can span millimeters to centimeters \cite{li2023efficient, turner22}. 
Using a uniform grid fine enough to resolve the melt pool is therefore computationally intractable, whereas general-purpose adaptive mesh refinement techniques incur significant overhead from frequent refinement, repartitioning, and inter-processor communication. These issues have spurred the development of GPU-accelerated multiresolution thermal solvers specialized for LPBF, including GO-MELT, HERMES, and other structured, nested moving-domain strategies \cite{aydin2026hermes, leonor2024go, li2023efficient}, as illustrated in \cref{fig:hermes_moving_domain}.
%, these methods  domains to separate localized near-laser thermal features from long-range heat diffusion.

While this spatial decomposition effectively mitigates the multiscale difficulty, the overall solution procedure remains inherently sequential in time. Along any given scan path, the temperature field at a given moment depends on the complete prior thermal history. Consequently, the wall-clock cost still grows with the entire simulated time horizon, making long scan paths and multi-layer builds expensive for design-space exploration and optimization.

\paragraph{Contributions}
The central observation driving this work is that, although LPBF thermal evolution is globally causal, the numerically relevant impact of sufficiently old scan-path segments diminishes with both time and spatial separation \cite{aydin2026hermes, yeung2020residual}. This motivates an alternative computational perspective: rather than enforcing a strictly serial progression, we decompose the path into contiguous segments and retain only the dominant, predecessor-induced thermal residues required to approximate each segment’s inherited state. This yields a sparse dependency pattern captured by a directed acyclic graph (DAG), in which nodes correspond to path segments and edges represent only thermally significant upstream interactions.

Within this framework, we introduce a segmented parallel-in-time formulation for LPBF thermal simulation. Our method substitutes the full step-by-step dependency chain with a sparse, segment-level DAG (constructed in \cref{seq:dag}), thereby exposing concurrency along the scan path and reducing the critical path to the depth of the graph. This, in turn, enables distributed execution across GPUs and compute nodes while bounding the approximation error with respect to the fully sequential reference. We detail the formulation and analysis of the method and empirically assess its behavior on multiple scan paths. For distributed-memory parallelism, we rely on the message passing interface (MPI) (see \cref{seq:partitioning}). We evaluate both accuracy~\cref{seq:accuracy} and strong scaling~\cref{seq:strong}. 

Although we employ a specific open-source solver~\cite{aydin2026hermes} for each time step, \method is solver-agnostic and can be integrated with any moving-domain solver. Additional related work is discussed in \cref{seq:related}. We discuss  limitations in \cref{seq:conclusions}.



% \subsection{Methodology}
% - Description of the main idea of parallel time using just a sequential path. \\
% - The definition of graph: what are nodes and edges\\
% - Graph construction and cycle detection\\
% - Graph partitioning for in parallel calculation\\
% - Load balancing (multiple nodes)\\
% - Single Node in-time parallelism\\
% - Summary of algorithm and how we select parameters (the edge definition as a function of the accuracy)\\
% - Error analysis \\
% - Parallel complexity \\

% \subsection{GPU/MPI implementation} 
% - Streams for in-parallel\\
% - MPI correction

% \subsection{Results}
% todo items besides writeup
% \begin{itemize}
% \item Figure out at what spatial resolution is the GPU underutilized
% \item Finalize settings; shapes, and layers, a laser paths. (Greenland, Texas, a couple of laser paths)
% \item Finalize accuracy goals and generate DAGS and partitioning to GPUs. 
% \item Decide on runs and scenarios, and number of GPUs

% %%%%
% \item Timing for uniform split vs our partitioning (using spiral).
% \item DAG demonstration (straight path and spiral path)
% \item Scalings
% \item Implement on texas path and island idea
% %%%%
% \end{itemize}

% Runs: For every dag and  two accuracy levels: low and high, do strong scaling from 1-GPU up to number of segments. 
% Can easily do strong scaling with multiple layers; 

% - Straight line -- to show convergence for different accuracy \\ 
% - Complex shapes, one layer linear -- to show to show different dags or different laser pathk\\
% - Complex shape linear multilayer \\
% - High-accuracy runs speedup (weak scaling)?\\
% - Very-low accuracy, underresolved speed up (kind of strong scaling). This will show issues with underutilization of single GPU: the work of a single time step it doesn't have enough blocks to keep a GPU busy. 
% - Maybe: Complex shape, one layer, nonlinear\\


% Speedup estimates compared to HERMES and open-foam from literature 


% \subsection{Related work}
% - OpenFOAM and ORNL paper \\
% - HERMES \\
% - Parallel in time for heat equation \\

% link with shapes https://sites.cc.gatech.edu/projects/large_models/blade.html
%% Greenland map



\section{Methodology}
\label{seq:methodology}



This section introduces the segmented parallel-in-time formulation \method\ proposed for the history-dependent LPBF thermal problem. The key idea is to substitute the conventional sequential dependency chain with a parallelizable two-stage approach: first, independent source-on simulations are performed; second, these are followed by parallel, sparse history corrections. By sparse, we mean that it is unnecessary to recompute the full scan path—only localized portions must be revisited. As will be shown, the computational cost of these corrections is governed by the desired target accuracy.  

We make the following assumptions:
\begin{enumerate*}
\item The heat source (laser) has constant parameters, such as power and scanning speed.
\item The thermal problem is linear with constant thermal conductivity.
\item Only a single laser is considered.
\end{enumerate*}
We discuss ways to relax these assumptions in \cref{seq:conclusions}.
%The first two assumptions are introduced for convenience, and the algorithm remains unchanged when they are relaxed.
%The extension to multiple lasers is beyond the scope of this paper.

\subsection{Segment-Wise Source-On/Source-Off Decomposition}
To demonstrate the core idea of \method, consider a laser that scans along a straight line $x \in [x_0,x_3]$ during the time interval $t \in [0,t_3]$. We divide this path into three equal segments, $s_1$, $s_2$, and $s_3$, corresponding to the time intervals $[0,t_1]$, $[t_1,t_2]$, and $[t_2,t_3]$, respectively, as illustrated in \cref{fig:three_segments}. 
The segments are such so that in during the first time interval the laser travels from left to right along $s_1$; during the second time interval the laser travels from left to right along $s_2$; and likewise for $s_3$. Let $u_1$, $u_2$, and $u_3$ global fields (defined in $[x_0,x_3]$) representing the thermal solutions associated with these three segments.

\begin{figure}[h]
    \centering
    \includegraphics[width=0.95\linewidth]{ThreeSegments.pdf}
    \caption{Straight-path example used to illustrate the segment-wise approach. The laser traverses three consecutive segments, $s_1$, $s_2$, and $s_3$, over the corresponding time intervals $[0,t_1]$, $[t_1,t_2]$, and $[t_2,t_3]$.}
    \label{fig:three_segments}
\end{figure}

In a standard sequential simulation, each segment inherits the terminal state of the previous one as its initial conditions. 
The PDEs below are defined in $[x_0,x_3]$
\begin{equation}
\begin{aligned}
s_1:\ &
\frac{\partial u_1}{\partial t} = \Delta u_1 + f_1,&\ t\in[0,t_1],
&\ u_1(0)=0, \\
 s_2:\ &
\frac{\partial u_2}{\partial t} = \Delta u_2 + f_2,&\ t\in[t_1,t_2],
&\ u_2(t_1)=u_1(t_1), \\
 s_3:\ &
\frac{\partial u_3}{\partial t} = \Delta u_3 + f_3,&\ t\in[t_2, t_3],
&\  u_3(t_2)=u_2(t_2), 
\end{aligned}
\label{eq:sequential_three_segments}
\end{equation}
where $f_i$ denotes the heat source during the $[t_{i-1},t_i]$. 
This creates a strict sequential dependency across the scan path.

The key idea in \method\ is to decompose each segment solution into an independent \emph{source-on} component and a \emph{source-off} history correction. 
For example, for segment $s_3$ we write
\begin{equation}
u_3 = u_3^* + v_3,
\label{eq:u3_decomp}
\end{equation}
where
\begin{equation}
\begin{aligned}
\frac{\partial u_3^*}{\partial t} &= \Delta u_3^* + f_3,
\qquad u_3^*(t_2)=0, \\
\frac{\partial v_3}{\partial t} &= \Delta v_3,
\qquad\qquad\ \, v_3(t_2)=u_2(t_2).
\end{aligned}
\label{eq:u3_star_v3}
\end{equation}
Here, $u_3^*$ represents the local contribution of the active source on $s_3$, while $v_3$ accounts for the thermal history inherited from earlier segments.

The central observation is that this inherited history need not always include the full upstream chain. If the segment length is sufficiently large, then the thermal influence of $s_1$ may decay below a prescribed tolerance by the time the laser reaches $s_3$. In that case, the contribution of $s_1$ to $u_3$ can be neglected, and the correction for $s_3$ may be approximated using only the immediate predecessor $s_2$. \cref{fig:segment_length_error} illustrates this effect for the three-segment straight-path example: as the segment length $L_{\mathrm{seg}}$ increases, the relative $L^2$ error from neglecting the contribution of $s_1$ at the start of $s_3$ decays rapidly.

\begin{figure}[h]
    \centering
    \includegraphics[width=0.95\linewidth]{Segment_L_vs_Error.pdf}
    \caption{Relative $L^2$ error at the start of $s_3$ versus segment length $L_{\mathrm{seg}}$ for the straight-path example in Fig.~\ref{fig:three_segments}. As the separation between $s_1$ and $s_3$ increases, the omitted contribution of $s_1$ decays rapidly, motivating the truncation of distant predecessors below a prescribed tolerance.}
    \label{fig:segment_length_error}
\end{figure}

This shortens the dependency chain substantially: the source-on solves $u_i^*$ are fully independent, and only a reduced set of correction solves must be propagated. In other words, shorter segments expose more concurrency but may require more history corrections, whereas longer segments reduce the number of relevant predecessors by increasing the spatiotemporal separation between segments.  This path-dependent truncation is related to near-far-field decomposition in N-body methods for heat potentials; 
it is the key mechanism that enables the high parallel efficiency of \method. The remainder of this section formalizes and builds on this idea.

\subsection{Segment Dependencies and Truncation Criteria}

We now generalize the three-segment example to a path divided into $N$ line segments $\{s_i\}_{i=1}^N$ corresponding to the movement of the laser, and the associated time intervals $[t_i,t_{i+1}]$. For each segment $s_r$, the total thermal solution is decomposed into a local source-on contribution and inherited source-off corrections:

\begin{equation}
u_r = u_r^* + \sum_{p<r} v_{p\to r}.
\label{eq:full_form}
\end{equation}

\paragraph{Source-on Simulation ($u^*$)} 
For each segment $s_r$, we define the \textbf{source-on} operator, $\mathcal{S}^{\mathrm{on}}_r(u_{\mathrm{init}})$, over the local time interval $[t_r,t_{r+1}]$. For the source-on solve, every segment starts from the reference state $u_{\mathrm{init}} = 0$:
\begin{equation}
u_r^* = \mathcal{S}^{\mathrm{on}}_r(0) \implies 
\begin{cases} 
\frac{\partial u^*}{\partial t} = \Delta u^* + f_r, \\
u^*(t_r)=0.
\end{cases}
\label{eq:source_on}
\end{equation}
Here, $f_r$ denotes the heat source restricted to segment $s_r$. Since these solves are initialized independently of all previous segments, the fields $\{u_r^*\}_{r=1}^N$ are mutually independent.

\paragraph{History Correction ($v$)} 
To account for heat accumulated from previous laser passes, a correction $v$ is applied. This is governed by the \textbf{source-off} operator, $\mathcal{S}^{\mathrm{off}}_r(u_0)$, which evolves an initial state $u_0$ across the time interval $[t_r,t_{r+1}]$ with the laser turned off:
\begin{equation}
v = \mathcal{S}^{\mathrm{off}}_r(u_0) \implies 
\begin{cases} 
\frac{\partial v}{\partial t} = \Delta v, \\
v(t_r) = u_0.
\end{cases}
\label{eq:source_off}
\end{equation}
The correction term $v_{p\to r}$ in \cref{eq:full_form} denotes the thermal tail generated on segment $s_p$ and propagated to the time interval of segment $s_r$. These corrections are computed recursively:
\begin{equation}
v_{p\to r} = 
\begin{cases} 
\mathcal{S}^{\mathrm{off}}_r\!\left(u_p^*(t_r)\right), & \text{if } p = r - 1, \\[0.5ex]
\mathcal{S}^{\mathrm{off}}_r\!\left(v_{p\to p+1}(t_r)\right), & \text{if } p+1 < r.
\end{cases}
\label{eq:correction_recursion}
\end{equation}
Thus, the contribution from the immediate predecessor is obtained by propagating its source-on field one segment forward, while the contribution from an older predecessor is obtained by repeatedly applying the source-off operator across intermediate segments. Because the governing problem is linear, these inherited corrections are additive, which justifies the superposition in \cref{eq:full_form}.


Although the thermal evolution is globally causal, the influence of a predecessor segment decays with elapsed time and spatial separation. We therefore retain a correction only if omitting it would induce a relative $L^2$ error exceeding a prescribed tolerance $e_{\mathrm{tol}}$. For a predecessor $s_i$ and target segment $s_j$, the exact retention criterion is
\begin{equation}
\max_{t \in (t_j,t_{j+1}]}
\frac{
\left\|u_j^{\mathrm{with}\,s_i}(t)-u_j^{\mathrm{without}\,s_i}(t)\right\|_{L^2}
}{
\left\|u_j^{\mathrm{with}\,s_i}(t)\right\|_{L^2}
}
> e_{\mathrm{tol}},
\label{eq:edge_criterion}
\end{equation}
where the norm is evaluated over a local region around the laser position at time $t$. This yields the truncated representation
\begin{equation}
u_r \approx u_r^* + \sum_{p \in \mathcal{P}_{e_{\mathrm{tol}}}(r)} v_{p\to r},
\label{eq:truncated_form}
\end{equation}
where $\mathcal{P}_{e_{\mathrm{tol}}}(r)$ denotes the set of predecessors whose omission would exceed the prescribed tolerance.

\subsection{Segment-Length Selection and Isotherm Approximation}
\label{seq:calibration_isotherm}
Here we explain how to determine when two segments are sufficiently separated in space and time so that their thermal interactions can be neglected. 
To assess whether an earlier segment can affect a later one, we require a simple rule for spatiotemporal connectivity. 
Our method relies on an influence radius: for a given elapsed time $\tau$, we introduce a radius $r(\tau)$ such that points located at distances greater than $r(\tau)$ from the source are regarded as thermally insignificant. The magnitude of this radius is controlled by a temperature threshold $\varepsilon$, which specifies the minimum residual thermal effect we still regard as relevant.

\paragraph{Parameter Selection}
We set he segment length $L_{\mathrm{seg}}$ and the threshold $\varepsilon$ using the straight-line experiment introduced in \cref{fig:segment_length_error}. 
We term this parameter selection as \emph{calibration} and the selected parameters as \emph{calibrated} parameters.
The path is divided into three equal segments of length $L$, and a corrected run in which $s_3$ uses only $s_2$ as predecessor is compared against the sequential baseline \cite{aydin2026hermes}. The observed error is the relative $L^2$ difference at the start of $s_3$, measured over a fixed evaluation region centered on the laser position. As shown in \cref{fig:segment_length_error}, this error decays rapidly with increasing $L$. We choose $L_{\mathrm{seg}}$ so that this error meets prescribed  tolerance.

The same experiment also records the maximum predecessor-induced temperature change in the target region. We denote this temperature change by $\varepsilon$. Intuitively, $\varepsilon$ specifies the smallest residual thermal effect that we still regard as significant when deciding whether two segments are connected. A larger $\varepsilon$ leads to a smaller influence radius, while a smaller $\varepsilon$ leads to a larger one. Since each value of $\varepsilon$ corresponds to an observed relative $L^2$ error, \cref{tab:calibration} provides a practical mapping from the user-prescribed error tolerance $e_{\mathrm{tol}}$ to the temperature threshold $\varepsilon$ used in the geometric connectivity test.

\begin{table}[h]
\centering
\caption{Parameter selection results for the straight-path experiment. For each segment length $L_{\mathrm{seg}}$, the table reports the maximum predecessor-induced temperature change $\varepsilon$ in the target segment's evaluation region together with the corresponding relative $L^2$ error.}
\label{tab:calibration}
\begin{tabular}{ccc}
\hline
$L_{\text{seg}}$ (mm) & $\varepsilon$ (K) & Rel.\ $L^2$ error \\
\hline
0.6 & $4.07\times10^{2}$ & $6.62\times10^{-3}$ \\
0.7 & $1.15\times10^{2}$ & $1.87\times10^{-3}$ \\
0.8 & $2.74\times10^{1}$ & $4.50\times10^{-4}$ \\
0.9 & $5.72\times10^{0}$ & $9.60\times10^{-5}$ \\
1.0 & $1.09\times10^{0}$ & $1.87\times10^{-5}$ \\
1.1 & $1.93\times10^{-1}$ & $3.41\times10^{-6}$ \\
1.2 & $3.24\times10^{-2}$ & $5.90\times10^{-7}$ \\
1.3 & $5.23\times10^{-3}$ & $9.79\times10^{-8}$ \\
1.4 & $8.19\times10^{-4}$ & $1.58\times10^{-8}$ \\
\hline
\end{tabular}
\end{table}

\paragraph{Isotherm approximation}
To construct the connectivity rule, we model the thermal remnant of a source point using the decay of the heat equation,
\begin{equation}
\Delta T(d,\tau) \propto \tau^{-3/2}\exp\!\left(-\frac{d^2}{4\alpha \tau}\right),
\label{eq:point_source_decay}
\end{equation}
where $d$ is the distance from the source, $\tau$ is the elapsed time, and $\alpha$ is the thermal diffusivity. For a prescribed threshold $\varepsilon$, this defines an influence radius $r(\tau)$ by
\begin{equation}
r(\tau) = \max\left\{d : \Delta T(d,\tau)=\varepsilon\right\}.
\label{eq:isotherm_radius}
\end{equation}
Equivalently, $r(\tau)$ is the largest distance at which the residual temperature remains at least $\varepsilon$ after time $\tau$.

As illustrated in \cref{fig:isotherm}, this influence region first expands and then contracts as the heat pulse diffuses and dissipates. In the next subsection, we use this radius to determine whether two segments are connected: if the target segment lies within the predecessor's influence radius at the relevant elapsed time, we retain a dependency; otherwise, we neglect it.

\begin{figure}[h]
    \centering
    \includegraphics[width=0.95\linewidth]{figs/Isotherm.pdf}
    \caption{spatiotemporal dependency evaluation via dynamic isotherm evolution. The thermal influence zone of source segment $s_i$ expands and then contracts over elapsed time. A dependency is retained if and only if the target segment intersects the active $\varepsilon$-isotherm at the corresponding time delay.}
    \label{fig:isotherm}
\end{figure}

In practice, $r(\tau)$ depends only on the material properties and the chosen threshold $\varepsilon$, and not on the particular scan geometry. We therefore precompute $r(\tau)$ once and store it as a look-up table for use in the DAG construction.











% \subsection{Path Discretization and Processor Assignment}
% We present the basic idea in \method using a simple example.
% In \cref{fig:discretization}, we show a simple laser path (essentially looking at the surface of the LBPF system from the top).
% First, we divide the path into $N$ line segments $\{s_i\}_{i=1}^N$ corresonding to the movement of the laser, and the associated time intervals $[t_i,t_{i+1}]$.  All the segments have the same length



% As illustrated in \cref{fig:discretization}, the simulation is decomposed into two stages. First, the physical scan path is divided into $N$ fine segments $s_i$, representing the smallest units of laser movement and associated time intervals $[t_i, t_{i+1}]$. Second, these segments are grouped into contiguous super-segments assigned to $M$ parallel processors (e.g., GPUs), denoted $P_0, P_1, \dots, P_{M-1}$.

% \begin{figure}[h]
%     \centering
% \includegraphics[width=\linewidth]{Discretization.pdf}    \caption{Discretization strategy: A) The physical raster path is discretized into fine segments $s_i$. B) Segments are grouped and assigned to processors.}
%     \label{fig:discretization}
% \end{figure}

% % \begin{figure}[h]
% %     \centering
% %     \fbox{\rule{0pt}{2in} \rule{\linewidth}{0pt}}
% %     \caption{Discretization strategy: (1) The physical raster path is discretized into fine segments $s_i$. (2) Segments are linearized and grouped into super-segments  for processor assignment. (3) Thermal influence $v$ from a previous processor $P_0$ crosses the boundary into $P_1$, necessitating a correction step.}
% %     \label{fig:discretization}
% % \end{figure}

% For this study, we utilize the linear outer model of HERMES and neglect radiative contributions (see \cref{fig:hermes_moving_domain}). Following the HERMES formulation, we define the non-dimensional temperature $u$ as:
% \begin{equation}
% u = \frac{T - T_s}{T_l - T_s},
% \label{eq:u_definition}
% \end{equation}
% where $T_s$ and $T_l$ are the solidus and liquidus temperatures, respectively. The resulting governing equation is:
% \begin{equation}
% \frac{\partial u}{\partial t} = \Delta u \qquad \text{in } \Omega.
% \label{eq:linear_heat}
% \end{equation}
% We refer the reader to the original HERMES work \cite{aydin2026hermes} for the detailed derivation of the scaling factors. Boundary conditions include the laser input $q_{\mathrm{surf}}^*$ and linear convection $q_{\mathrm{conv}}$ on the top surface $\partial\Omega_{\mathrm{top}}$.

\subsection{Dependency-DAG Construction}
\label{seq:dag}

We now use the selected segment length $L_{\mathrm{seg}}$ and the precomputed influence radius $r(\tau)$ to construct the dependency DAG. Each node of the graph represents one scan segment, and a directed edge is added only when the thermal remnant of an earlier segment remains significant at a later one.

\paragraph{Nodes and edges}
Let $\mathcal{G} = (\mathcal{V},\mathcal{E})$ denote the dependency graph, where each node $s_i \in \mathcal{V}$ is a scan segment of length $L_{\mathrm{seg}}$. Since thermal influence propagates only forward in time, we only consider edges of the form $(s_i,s_j)$ with $i<j$.

To determine whether such an edge should be kept, we compare the spatiotemporal separation between the two segments against the influence radius. For a segment pair $(s_i,s_j)$, let
\begin{equation}
\tau_{ij} = t_j - t_{i+1}
\label{eq:tauij}
\end{equation}
denote the elapsed time from the end of segment $s_i$ to the start of segment $s_j$. We also define
\begin{equation}
d_{ij} = \min_{x \in s_i,\; y \in s_j} \|x-y\|,
\label{eq:dij}
\end{equation}
as the minimum Euclidean distance between the two segments. This gives a conservative connectivity test: if even the closest points of the two segments lie outside the influence radius after the elapsed time $\tau_{ij}$, then the segments are treated as independent.

We retain the edge $(s_i,s_j)\in\mathcal{E}$ if and only if
\begin{equation}
d_{ij} \le r(\tau_{ij}).
\label{eq:dag_check}
\end{equation}
In words, a future segment depends on a predecessor precisely when the predecessor's thermal influence region still reaches the future segment after the corresponding elapsed time. Otherwise, the predecessor is omitted from the dependency graph. Therefore, the resulting DAG directly determines the correction workload: each incoming edge to a segment corresponds to a predecessor contribution that must be retained and propagated through a source-off correction. Hence, the sparsity of the DAG controls how many correction solves are required.

\paragraph{Interpretation}
The threshold $\varepsilon$ chosen in the  step controls the size of the radius $r(\tau)$. A larger value of $\varepsilon$ yields a smaller influence radius and therefore a sparser graph, whereas a smaller value of $\varepsilon$ yields a larger influence radius and a denser graph. Through the calibration table, this allows the user to prescribe an error tolerance $e_{\mathrm{tol}}$ and obtain a corresponding temperature threshold $\varepsilon$, which then determines the graph sparsity through \cref{eq:dag_check}.

\paragraph{Complexity}
Since $r(\tau)$ vanishes beyond a finite time horizon, each segment can influence only a bounded number of future segments. If each segment affects at most $M$ successors, the DAG construction cost is $O(N\cdot M)$.

\subsection{Error Analysis and Stability}
\label{seq:error_analysis}

The accuracy of \method\ is governed by the omission of predecessor contributions that are not retained in the DAG. For the linear model considered here, the resulting error can be interpreted as the perturbation introduced by removing those contributions from the exact sequential solution.

\paragraph{Global error decomposition}
Let $u_j^{\mathrm{exact}}$ denote the temperature at segment $s_j$ when all predecessor contributions are retained, and let $u_j$ denote the truncated approximation in \cref{eq:truncated_form}. The error at segment $s_j$ is then
\begin{equation}
E_j = u_j^{\mathrm{exact}} - u_j
    = \sum_{i \in \mathcal{P}_{\mathrm{ignored}}(j)} v_{i \to j},
\label{eq:global_error_sum}
\end{equation}
where $\mathcal{P}_{\mathrm{ignored}}(j)$ is the set of predecessors $i<j$ whose edges to $s_j$ were omitted from the DAG, and $v_{i\to j}$ is the corresponding thermal tail propagated to $s_j$ via \cref{eq:correction_recursion}. Thus, the truncation error is exactly the sum of the discarded thermal histories.

\paragraph{Local truncation and physical decay}
By construction of the dependency check in \cref{eq:dag_check}, every omitted edge $(s_i,s_j)$ corresponds to a predecessor contribution whose residual temperature at the target is below the calibrated threshold $\varepsilon$. From the calibration procedure in \cref{tab:calibration}, this threshold corresponds to a relative $L^2$ omission error of approximately $e_\varepsilon$. Therefore, if at most $M$ predecessor contributions are omitted for any target segment, the segmentwise error satisfies the bound
\begin{equation}
\|E_j\|_{L^2} \le M\, e_\varepsilon.
\label{eq:global_error_bound}
\end{equation}
This bound is conservative. In practice, predecessors that are far away in space or time contribute errors far smaller than the calibrated threshold, so the effective error is usually dominated by segments near the truncation boundary, as also confirmed by Tables~\ref{tab:acc_1e4} and~\ref{tab:acc_1e7}.

\paragraph{Contractive stability}
The source-off operator $\mathcal{S}^{\mathrm{off}}$ is contractive in the $L^2$ sense:
\begin{equation}
\|\mathcal{S}^{\mathrm{off}}(u)\|_{L^2} \le \|u\|_{L^2}.
\label{eq:stability_contraction}
\end{equation}
As a result, approximation errors introduced by truncating earlier segments do not amplify as they propagate forward. Instead, they decay under diffusion, which ensures stability of the truncated evolution over the full scan path.


% \subsection{Error Analysis and Stability}
% \label{seq:error_analysis}

% The accuracy of the \method is governed by the
% truncation of the dependency DAG. For a linear model, we can formalize the error by treating the truncation as a perturbation of the exact solution.
% \paragraph{Global Error Decomposition}
% Let $u_{j}^{\text{exact}}$ be the temperature at segment $s_j$ if all
% previous history were included, and let $u_{j}$ be the approximated solution using \cref{eq:truncated_form}. The global error $E_j$ is:
% \begin{equation}
% E_j = u_{j}^{\text{exact}} - u_{j} = \sum_{i \in \mathcal{P}_{\text{ignored}}(j)} v_{i \to j},
% \label{eq:global_error_sum}
% \end{equation}
% where $\mathcal{P}_{\text{ignored}}(j)$ contains all segments $i < j$
% without a retained edge to $s_j$, and $v_{i \to j}$ is the thermal tail of
% $s_i$ propagated to $s_j$ via \cref{eq:correction_recursion}.

% \paragraph{Local Truncation and Physical Decay}
% By construction of the dependency check in \cref{eq:dag_check}, every
% truncated edge $(s_i \to s_j)$ corresponds to a thermal contribution whose
% top-surface temperature at the target ROI is below $\varepsilon$. From the
% calibration experiment in Section~\ref{sec:dag}, this threshold
% $\varepsilon$ corresponds to a per-truncation relative $L^2$ error of
% $e_\varepsilon$. The global error is therefore bounded by:
% \begin{equation}
% E_{\mathrm{global}} \leq M \cdot e_\varepsilon,
% \label{eq:global_error_bound}
% \end{equation}
% where $M$ is the maximum number of predecessors any segment can have,
% determined by the isotherm decay rate and threshold $\varepsilon$ (\cref{fig:isothermal}). In practice, this bound is
% highly conservative: truncated segments far from the target in space or
% time contribute errors several orders of magnitude below $e_\varepsilon$,
% and the effective error is dominated only by segments near the truncation
% boundary, as confirmed by Tables~\ref{tab:acc_1e4} and~\ref{tab:acc_1e7}.

% \paragraph{Contractive Stability}
% The source-off operator $\mathcal{S}^{\text{off}}$ is a contractive mapping
% in the $L^2$ sense:
% \begin{equation}
% \|\mathcal{S}^{\text{off}}(u)\| \le \|u\|.
% \label{eq:stability_contraction}
% \end{equation}
% This ensures that approximation errors introduced in early processors do not
% amplify as they propagate forward in time; instead, the error energy
% naturally decays due to thermal diffusion and convective losses, guaranteeing
% numerical stability across the entire simulation horizon.




% \paragraph{Sparse Truncation}
% Although LPBF thermal evolution is globally causal, the influence of a
% source segment decays rapidly with both elapsed time $\tau$ and spatial
% distance $d$. Motivated by the point-source solution to the heat equation,
% the temperature rise at distance $d$ after time $\tau$ scales as:
% \begin{equation}
% \label{eq:decay_rate}
% \Delta T(d,\tau) \propto \tau^{-3/2} \exp\left(-\frac{d^2}{4\alpha\tau}\right),
% \end{equation}
% where $\alpha$ is the thermal diffusivity. For a prescribed temperature
% threshold $\varepsilon$, this decay defines a time-evolving isotherm on the
% top surface whose radius first expands and then contracts as the thermal
% pulse diffuses and dissipates, as depicted in Fig.~\ref{fig:isothermal}.

% We exploit this decay to sparsify the dependency structure: the correction
% $v_{p \to r}$ in \cref{eq:full_form} is retained only if the estimated
% thermal influence of $P_p$ at the boundary of $P_r$ exceeds $\varepsilon$.
% This turns the otherwise dense sequential chain into a sparse DAG,
% whose construction is detailed in Section~\ref{sec:dag}.
% The truncated solution is:
% \begin{equation}
% \label{eq:truncated_form}
% u_r \approx u^*_r + \sum_{p \in \mathcal{P}_\varepsilon(r)} v_{p \to r},
% \end{equation}
% where $\mathcal{P}_{\varepsilon}(r) = \{ p < r : \mathrm{Influence}(P_p
% \to P_r) > \varepsilon \}$ denotes the set of thermally significant
% predecessors.


 

% \subsection{DAG Construction}
% \label{seq:dag}

% % We build a DAG in which each node corresponds to a segment $s_i$. A directed edge $s_i \to s_j$ ($i < j$) is kept only if neglecting the thermal contribution of $s_i$ to $s_j$ would result in a relative $L^2$ error greater than $e_{\mathrm{tol}}$. \emph{Since thermal effects propagate only forward in time, the resulting graph is, by definition, acyclic.}

% \paragraph{Calibration}
% We first calibrate the segment length $L_{\mathrm{seg}}$ via a controlled
% straight-line experiment. For each laser position, we define an
% \emph{evaluation region} as a $1\times1\times1$~mm box whose top face is
% centered on the laser position and extends 1~mm downward in the
% $z$-direction. The path is divided into three equal segments of length $L$;
% a corrected run in which $s_3$ uses only $s_2$ as predecessor is compared
% against a fully sequential baseline. The observed error is the relative
% $L^2$ difference over the evaluation region at the start of $s_3$. As shown
% in Fig.~\ref{fig:calibration_error}, this error decays exponentially with
% $L$. Once the $L_{\mathrm{seg}}$ value corresponding to the target accuracy is selected, it is applied to all trajectories, including non-straight paths. 

% \begin{figure}[h]
%     \centering
%     \includegraphics[width=\linewidth]{straight_line_error.pdf}
%     \caption{Relative $L^2$ error at the start of $s_3$ as a function of
%     segment length $L$ for a straight-line path. The exponential decay
%     confirms that beyond a sufficient separation, the truncated contribution
%     of $s_1$ drops below the numerical noise floor.}
%     \label{fig:calibration_error}
% \end{figure}





% \paragraph{Graph: Node and Edge}
% The thermal dependency structure between segments is represented as a graph $\mathcal{G} =
% (\mathcal{V}, \mathcal{E})$, where each node $s_i \in \mathcal{V}$ is a
% scan segment of length $L_{\mathrm{seg}}$. A directed edge $(s_i, s_j) \in \mathcal{E}$ ($i < j$) is retained if and
% only if:
% \begin{equation}
% \max_{t \in (t_j, t_{j+1}]} \frac{\|u_j^{\mathrm{with}\,s_i}(t) - u_j^{\mathrm{without}\,s_i}(t)\|_{L^2}}
% {\|u_j^{\mathrm{with}\,s_i}(t)\|_{L^2}} > e_{\mathrm{tol}},
% \label{eq:edge_criterion}
% \end{equation}
% where the $L^2$ norm is evaluated over the evaluation region centered on
% the laser position at time $t$. Because thermal influence
% only travels forward in time, the graph is acyclic by construction, yielding
% a DAG.


% % \paragraph{Path awareness}
% % By construction, the calibration ensures that on a straight-line path each
% % segment depends only on its immediate predecessor, yielding a sparse
% % near-chain structure. In contrast, complex trajectories such as raster or
% % spiral paths, where the laser repeatedly revisits nearby spatial regions,
% % generate denser dependency structures with longer-range edges, as illustrated
% % in \cref{fig:geometry_dags}.
% % This path-aware sparsity is critical for correctness. 


% % \begin{figure}[h]
% %     \centering
% %     \includegraphics[width=\linewidth]{DAG_three_cases.pdf}
% %     \caption{Dependency structures for different scan paths. A straight-line path (left) reduces to a simple chain. In contrast, raster (center) and spiral (right) paths generate denser dependencies where a single segment may inherit history from multiple spatially proximate predecessors.}
% %     \label{fig:geometry_dags}
% % \end{figure}









% \begin{table}[h]
% \centering
% \caption{Calibration results: maximum top-surface temperature $\varepsilon$
% in the ROI and relative $L^2$ error as a function of segment length $L$.}
% \label{tab:calibration}
% \begin{tabular}{ccc}
% \hline
% $L$ (mm) & $\varepsilon$ (K) & Rel $L^2$ error \\
% \hline
% 0.6 & $4.07\times10^{2}$ & $6.62\times10^{-3}$ \\
% 0.7 & $1.15\times10^{2}$ & $1.87\times10^{-3}$ \\
% 0.8 & $2.74\times10^{1}$ & $4.50\times10^{-4}$ \\
% 0.9 & $5.72\times10^{0}$ & $9.60\times10^{-5}$ \\
% 1.0 & $1.09\times10^{0}$ & $1.87\times10^{-5}$ \\
% 1.1 & $1.93\times10^{-1}$ & $3.41\times10^{-6}$ \\
% 1.2 & $3.24\times10^{-2}$ & $5.90\times10^{-7}$ \\
% 1.3 & $5.23\times10^{-3}$ & $9.79\times10^{-8}$ \\
% 1.4 & $8.19\times10^{-4}$ & $1.58\times10^{-8}$ \\
% \hline
% \end{tabular}
% \end{table}




% \paragraph{Approximate Model: Isotherm Method}
% Directly evaluating \cref{eq:edge_criterion} for every segment pair at
% runtime is infeasible. Instead, we exploit the calibration experiment, which
% simultaneously records the relative $L^2$ error and the maximum temperature
% on the top surface of the evaluation region for each segment length $L$.
% This establishes a direct correspondence: for a target tolerance
% $e_{\mathrm{tol}}$, there exists a corresponding temperature threshold
% $\varepsilon$ such that the relative $L^2$ error exceeds $e_{\mathrm{tol}}$
% if and only if the maximum top-surface temperature induced by $s_i$ within
% the evaluation region of $s_j$ exceeds $\varepsilon$. We can therefore
% replace \cref{eq:edge_criterion} with an equivalent point temperature check:
% retain edge $(s_i, s_j)$ if and only if the maximum top-surface temperature
% induced by $s_i$ in the evaluation region of $s_j$ exceeds $\varepsilon$.

% For any segment pair $(s_i, s_j)$, we define the time delay $\tau_{ij} =
% t_j - t_i$ as the elapsed time between the two segments, and the spatial
% distance \begin{equation}
% d_{ij} = \min_{p \in s_i} \min_{t \in (t_j, t_{j+1}]} \mathrm{dist}(p,\, \mathcal{R}_j(t)),
% \label{eq:dij}
% \end{equation}
% where $\mathcal{R}_j(t)$ denotes the evaluation region centered on the
% laser position at time $t$.   To enable this test, we model each source segment as a series of point
% heat sources distributed along the scan path. For a point source, the
% temperature rise at distance $d$ after elapsed time $\tau$ is given by:
% \begin{equation}
% \label{eq:decay_rate}
% \Delta u(d,\tau) \propto \tau^{-3/2} \exp\left(-\frac{d^2}{4\alpha\tau}\right),
% \end{equation}
% where $\alpha$ is the thermal diffusivity. Given $d_{ij}$ and $\tau_{ij}$
% for a segment pair $(s_i, s_j)$, this expression allows us to estimate the
% temperature rise induced by $s_i$ within the evaluation region of $s_j$.


% Rather than evaluating $\Delta u$ for every $(d_{ij}, \tau_{ij})$ pair
% explicitly, we adopt an isotherm-based approach. For a given $\varepsilon$, the isotherm radius $r(\tau)$ is defined as the
% maximum distance on the top surface at which the temperature rise equals
% $\varepsilon$:
% \begin{equation}
% r(\tau) = \max \{ d : \Delta u(d, \tau) = \varepsilon \}.
% \label{eq:isotherm_radius}
% \end{equation}
% Since the laser deposits energy at the top surface, the maximum temperature
% rise occurs there, and it suffices to track the isotherm on the top surface
% alone. As illustrated in \cref{fig:isothermal}, a dependency edge
% $s_i \to s_j$ is retained if and only if the evaluation regions of $s_j$
% intersect the $\varepsilon$-isotherm of $s_i$ at time delay $\tau_{ij}$.




% \begin{figure}[h]
%     \centering
%     \includegraphics[width=\linewidth]{thermal_level.pdf}
%     \caption{spatiotemporal dependency evaluation via dynamic isotherm evolution. The thermal influence zone of source segment $s_i$ (the laser location at during $(t_i,t_{i+1}]$), expands and shrinks over the elapsed time $t_i+\delta t$. A dependency edge is retained in the DAG if and only if a future target segment's region of interest intersects the active isotherm at their specific time delay. As illustrated, target $s_{i+3}$ intersects the isotherm (establishing a thermal dependency with $s_i$), while $s_{i+4}$ resides safely outside (requiring no history correction from $s_i$).}
%     \label{fig:isothermal}
% \end{figure}

% \paragraph{Dependency Check}
% Since $r(\tau)$ depends only on the material properties through $\alpha$ and
% the threshold $\varepsilon$, and not on the specific scan path geometry, we
% precompute $r(\tau)$ into a look-up table built once per material
% configuration. For any segment pair $(s_i, s_j)$, $(s_i, s_j) \in \mathcal{E}$ if and only if:
% \begin{equation}
% d_{ij} \leq r(\tau_{ij}),
% \label{eq:dag_check}
% \end{equation} with edge weight $|j - i|$.
% respectively. Since $r(\tau) = 0$ beyond a finite time, each $s_i$ affects
% at most $M$ successors, giving a DAG construction complexity of $O(N \cdot M)$.






% % \begin{figure}[h]
% %     \centering
% %     \includegraphics[width=\linewidth]{detection.pdf}
% %     \caption{Coupled-area detection under different correction-depth budgets ($K$). Top: With a budget of $K=1$, the boundary (dashed line) is admissible because no edges overreach the limit. Bottom: With a budget of $K=2$, the boundary can accommodate denser dependencies without triggering a required merge of segments.}
% %     \label{fig:cut_detection}
% % \end{figure}



% % \begin{figure}
% %     \centering
% %     \includegraphics[width=\linewidth]{figs/Grouping.pdf}
% %     \caption{Super-segment extraction under a correction budget of $K=2$. \textbf{Top:} Each boundary is evaluated for its cut depth, determined by the maximum span of crossing dependency edges (colored arrows). \textbf{Bottom:} Partition cuts (dashed lines) are placed exclusively at boundaries where the cut depth satisfies $d_b \le 2$. The enclosed segments are fused into contiguous super-segments ($ss_1, \dots, ss_4$), safely encapsulating dense topological dependencies internally while leaving only shallow connections across block boundaries.}
% %     \label{fig:grouping}
% % \end{figure}


\subsection{DAG partitioning}
\label{seq:partitioning}

The preceding subsections define which predecessor contributions must be retained through the dependency DAG. We now address the computational question of how to assign this work to processors. Given $N$ segments and $M$ processors, we seek a partition of the scan path into $M$ contiguous set of segments, each assigned to one processor. Within a processor, segments are integrated sequentially, so all internal dependencies are resolved automatically. Across processor boundaries, however, retained DAG edges require explicit source-off corrections. The goal of the partitioning is therefore to balance the source-on workload while minimizing the correction overhead induced by cross-boundary dependencies.


\paragraph{Cut Depth}
For any contiguous range of segments $[i, j]$, we define the cut depth
$c(i,j)$ as the maximum reach of any dependency edge originating within
$[i,j]$ and crossing its right boundary:
\begin{equation}
c(i,j) = \max \{ k - j : (s_p, s_k) \in \mathcal{E},\ i \leq p \leq j,\ k > j \}.
\label{eq:cut_depth_def}
\end{equation}
In other words, $c(i,j)$ measures how far the thermal history within a processor must be propagated beyond its right boundary. Since segments within $[i,j]$ are solved sequentially, only edges crossing the right boundary require explicit source-off propagation. This means that if $[i,j]$ is assigned to a single processor, it must apply the source-off operator
$c(i,j)$ times after its source-on simulation to propagate its thermal history to all dependent successor segments, as illustrated in \cref{fig:cut_depth}.

% Since segments within $[i,j]$ are solved sequentially, the thermal
% history of any segment $s_p$ is already incorporated into the state of
% $s_{p+1}$, and ultimately into $s_j$. Only edges crossing the right
%boundary require explicit source-off propagation. 


\begin{figure}[h]
    \centering
    \includegraphics[width=\linewidth]{figs/cut_depth.pdf}
\caption{Illustration of cut depth for a partition into three processors.
The cut depth $c(1,5) = 4$ indicates that the set of segments $[1,5]$
must apply the source-off operator 4 times to propagate its thermal history
to all dependent successor segments. Similarly, $c(6,10) = 3$ for the
second set of segments.}
    \label{fig:cut_depth}
\end{figure}

\paragraph{Workload Model}
The computational cost of segments $[i, j]$ inside a single processor consists of the source-on simulation over all assigned segments and the source-off propagations at the exit boundary:
\begin{equation}
\mathrm{cost}(i,j) = (j-i+1) \cdot w_{\mathrm{source}} + c(i,j) \cdot w_{\mathrm{corr}},
\label{eq:cost}
\end{equation}
where $w_{\mathrm{source}}$ and $w_{\mathrm{corr}}$ are the wall-clock cost of a single source-on and source-off solve over one segment, respectively. Since all segments have the same length $L_{\mathrm{seg}}$, the per-segment costs $w_{\mathrm{source}}$ and $w_{\mathrm{corr}}$ are approximately constant across the scan path.

\paragraph{Uniform Partitioning}
Regardless of how segments are assigned to processors, the DAG is required: each processor must consult it to determine which source-off corrections to compute and where to send them. The simplest assignment is to divide the $N$ segments uniformly into $M$ equal-sized sets, assigning each to one processor. This guarantees balanced source-on workload, as each processor handles $\lfloor N/M \rfloor$ segments. However, uniform partitioning ignores the correction cost $c(i,j) \cdot w_{\mathrm{corr}}$ in \cref{eq:cost}: partition boundaries may fall in regions of dense thermal dependency, resulting in large cut depths and high correction overhead that degrades parallel efficiency.

\paragraph{Optimal Partitioning}
To jointly minimize source-on and correction costs, we seek a partition that
minimizes the bottleneck load across all processors. Let $DP[p,j]$ denote
the minimum bottleneck load when optimally partitioning the first $j$
segments across $p$ processors. The recurrence is:
\begin{equation}
DP[p,j] = \min_{i \leq j} \Big( \max \big( DP[p-1,i-1],\, \mathrm{cost}(i,j) \big) \Big),
\label{eq:dp_recurrence}
\end{equation}
with base case $DP[1,j] = \mathrm{cost}(1,j)$. The optimal partition is
recovered by backtracking through the DP table. In this way, the DAG not only identifies which predecessor corrections are required, but also drives how the scan path should be partitioned for efficient parallel execution.



\subsection{Parallel algorithms}
\label{seq:parallel_execution}

Once the partition is fixed, the retained dependency structure determines the execution schedule. Each processor advances its assigned contiguous sets of segments sequentially with the source active, thereby resolving all internal dependencies locally. After completing its source-on integration, each processor computes the retained source-off corrections for edges that cross its right boundary and sends the results to the corresponding downstream processors. The downstream processors then accumulate these received corrections by superposition, without performing any additional PDE solves. This realizes the truncated representation in \cref{eq:truncated_form}, with corrections propagated according to \cref{eq:correction_recursion}.

As illustrated in \cref{fig:thermal_influence}, the final solution on a processor is obtained by superposing its local source-on field with the correction fields received from upstream processors.

\begin{figure}[h]
    \centering
    \includegraphics[width=\linewidth]{Thermal_Influence.png} 
    \caption{Parallel execution across a processor boundary. While processor $P_1$ advances its locally assigned segments with the source active, it also receives the retained thermal history propagated from processor $P_0$. Only DAG edges crossing processor boundaries require such explicit source-off correction transfers.}
    \label{fig:thermal_influence}
\end{figure}


The execution proceeds in two stages. In the first stage, all processors concurrently perform their local source-on integrations. In the second stage, each processor propagates its thermal history forward by repeatedly applying the source-off operator, once per retained outgoing DAG edge, and sends each result to the corresponding downstream processor. The downstream processor accumulates the received fields by addition. Because these corrections are required only for sparse cross-boundary dependencies, the computation and communication overhead remain limited.

Moreover, correction propagation does not need to wait for the completion of all processors. As soon as a processor finishes its source-on integration, it can begin computing and sending its outgoing corrections while other processors are still in their source-on phase. This creates a pipelined schedule in which correction work overlaps with ongoing source-on execution on other processors.

\cref{fig:parallel_timing} illustrates this schedule for $M=4$. The source-on phase is fully concurrent, while each sender's correction propagation forms a short trailing tail whose length equals its cut depth. The processor with the largest cut depth determines the overall completion time. This is the key reason the method achieves high parallel efficiency despite the history dependence of the original sequential problem.

\cref{alg:pipeline} summarizes the complete workflow, from segment discretization and DAG construction to partitioning and parallel propagation of retained corrections.


\begin{algorithm}
\caption{Parallel-in-Time LPBF Thermal Simulation}
\label{alg:pipeline}
\begin{algorithmic}[1]
\Require Scan path, $M$, $r(\tau)$, $L_{\mathrm{seg}}$

\State Discretize scan path into segments $\{s_1, \dots, s_N\}$ of length $L_{\mathrm{seg}}$
\For{each segment $s_i$}
    \For{$j = i+1, i+2, \dots$ \textbf{until} $r(\tau_{ij}) = 0$}
        \If{$d_{ij} \leq r(\tau_{ij})$}
            \State Add directed edge $(s_i, s_j)$ to $\mathcal{E}$
        \EndIf
    \EndFor
\EndFor
\State Precompute $c(i,j)$ for all valid ranges $[i,j]$
\State Solve DP to obtain optimal partition of $\{s_1,\dots,s_N\}$ into $M$ sets

\For{each processor $P_r$ assigned to segments $[a_r, b_r]$ \textbf{in parallel}} \Comment{Stage 1}
    \State $u^*_r \gets \mathcal{S}^{\mathrm{on}}(0)$
    \State $v \gets u^*_r(t_{b_r+1})$ \Comment{Terminal state at end of assigned interval}
    \For{$k = 1, \dots, c(a_r, b_r)$} \Comment{Cut depth of $[a_r, b_r]$}
        \State $v \gets \mathcal{S}^{\mathrm{off}}(v)$
        \State Send $v$ to the processor owning segment $s_{b_r+k}$
    \EndFor
\EndFor

\For{each processor $P_r$ \textbf{in parallel}} \Comment{Stage 2: accumulation}
    \State Receive all corrections $v_{p \to r}$ from predecessor processors
    \State $u_r \gets u^*_r + \sum_{p} v_{p \to r}$
\EndFor

\end{algorithmic}
\end{algorithm}






 
\subsection{Parallel Complexity and Scalability}
\label{seq:parallel_complexity}

The parallel cost of \method\ is determined by two ingredients: the concurrent source-on work distributed across processors, and the retained source-off corrections induced by DAG edges that cross processor boundaries. Let $N$ be the total number of segments and $M$ the number of processors. Let $w_{\mathrm{source}}$ denote the cost of advancing one segment with the source active, and let $w_{\mathrm{corr}}$ denote the cost of one source-off correction step. \cref{fig:parallel_timing} illustrates this execution pattern: the source-on phase is fully concurrent across processors, while the retained correction phase forms a much shorter trailing tail.

\begin{figure}[h]
    \centering
    \includegraphics[width=\linewidth]{figs/parallel_complexity.pdf}
\caption{Parallel execution timeline for $M=4$ processors without truncation. Each processor performs its source-on simulation concurrently, then propagates source-off corrections to downstream processors. The final solution on each processor (right) is the superposition of its source-on field and all received corrections.}
    \label{fig:parallel_timing}
\end{figure}


\paragraph{Theoretical complexity}
In a fully sequential simulation, the cost scales as
\begin{equation}
\mathcal{C}_{\mathrm{seq}} = N\,w_{\mathrm{source}},
\label{eq:complexity_seq}
\end{equation}
since all $N$ segments are advanced one after another with the source active.
In \method, the source-on work is distributed across $M$ processors, while each processor additionally performs source-off propagation steps determined by its cut depth. The wall-clock time is governed by the most heavily loaded processor:
\begin{equation}
\mathcal{C}_{\mathrm{para}} = \underbrace{\frac{N}{M}\, w_{\mathrm{source}}}_{\text{Concurrent phase}} + \underbrace{\max_{r}\, c(a_r,b_r) \cdot w_{\mathrm{corr}}}_{\text{Correction phase}},
\label{eq:complexity}
\end{equation}
where $c(a_r,b_r)$ is the cut depth of processor $P_r$ as defined in \cref{eq:cut_depth_def}. The first term represents the ideal load-balanced source-on phase, while the second term represents the longest sender-side correction tail across all processors.

Because source-off propagation is substantially cheaper than source-on advancement, the correction term is typically much smaller than the concurrent source-on term. Moreover, the partitioning strategy in \cref{seq:partitioning} explicitly avoids placing processor boundaries in regions of dense thermal dependency, which keeps the cut depths small. As a result, \method achieves near-ideal speedup.
\section{Empirical results}\label{seq:results}
\subsection{System Setup}
All experiments were performed on the Vista cluster at the Texas Advanced
Computing Center (TACC), using NVIDIA GH200 Grace-Hopper Superchip nodes,
each equipped with a single H200 GPU with 96~GiB of HBM3 memory.
Nodes are interconnected via NVIDIA InfiniBand NDR at 400~Gb/s.
Inter-node communication is handled through MPI.
In the MPS experiments (\cref{seq:strong}), two ranks share one GPU.

\cref{fig:scan_paths} shows the four scan paths considered in this
study: a Bull-shaped raster path, a Texas-shaped raster path, a synthetic
hybrid spiral--raster path, and a Hilbert-curve path.
For the Bull and Texas cases, the scan paths are constructed from
image-based patterns with a horizontal extent of 50~mm and a scan-line
spacing of 200~\si{\micro\meter}.
The synthetic hybrid path is constructed from a repeated unit consisting of
an inward spiral of size $8\times8$~mm followed by an attached raster region
of size $8\times4$~mm; this unit is repeated five times.
The Hilbert path occupies a $40\times40$~mm square and uses a Hilbert curve
of order~7.

The key thermal and laser parameters are summarized in
\cref{tab:thermal_params}.
All experiments use a time step of 10~\si{\micro\second}.
Unless otherwise stated, the top-surface temperature threshold is
$\varepsilon = 0.01$~K for the DAG construction.
The moving computation-domain size, grid spacing, and resulting grid
resolution are summarized in \cref{tab:domain_setup}.
The domain listed in \cref{tab:domain_setup} refers to the moving
computation domain that follows the laser, not the full extent of the
printed part.

\begin{table}[h]
\centering
\caption{Thermal and laser parameters.}
\label{tab:thermal_params}
\begin{tabular}{lll}
\hline
Parameter              & Symbol      & Value \\
\hline
Laser power            & $Q$         & 200~W \\
Scan speed             & $v_s$       & 1.0~m/s \\
Beam half-width        & $r_b$       & 25~\si{\micro\meter} \\
Absorptivity           & $A$         & 0.55 \\
Thermal conductivity   & $k$         & 13.5~W/(m$\cdot$K) \\
Density                & $\rho$      & 7950~kg/m$^3$ \\
Heat capacity          & $c_p$       & 470~J/(kg$\cdot$K) \\
Ambient temperature    & $T_0$       & 298.15~K \\
\hline
\end{tabular}
\end{table}

\begin{table}[h]
\centering
\caption{Moving computation-domain and grid settings.}
\label{tab:domain_setup}
\begin{tabular}{ccc}
\hline
Grid spacing $h$ (\si{\micro\meter}) & Domain size (mm) & Grid points \\
\hline
18 & $4.8\times4.8\times2.4$ & $267\times267\times134$ \\
30 & $4.8\times4.8\times2.4$ & $161\times161\times81$ \\
40 & $4.8\times4.8\times2.4$ & $120\times120\times60$ \\

\hline
\end{tabular}
\end{table}

\begin{figure}[h]
    \centering
    \includegraphics[width=\linewidth]{scan_path.png}
    \caption{Scan paths used in the experiments: Bull, Texas, spiral-raster (one unit), and Hilbert.}
    \label{fig:scan_paths}
\end{figure}

\subsection{Accuracy Study}\label{seq:accuracy}
We evaluate whether a calibration obtained from a straight scan path
transfers to more complex scan paths.
For a prescribed target accuracy, we perform a straight-path calibration
to determine the segment length $L_{\mathrm{seg}}$ and temperature threshold
$\varepsilon$, following the procedure described in
\cref{seq:calibration_isotherm}. These calibrated values are then applied to
three representative scan paths: a straight line, the Bull path, and the
hybrid spiral--raster path.

The transient temperature field is sampled every 25 time steps throughout
each run. The observed error is defined as the maximum relative $L^2$ error
over all sampled snapshots; we also record the global maximum cut depth
(see \cref{seq:partitioning}) of the corresponding DAG.

Tables~\ref{tab:acc_1e4} and~\ref{tab:acc_1e7} summarize results for two
calibrated accuracy levels. Across both configurations, the straight-line
error meets the target by construction, and the errors on the more complex
paths remain close to the straight-line baseline in most cases.
The hybrid spiral--raster path under the $10^{-4}$ configuration
(\cref{tab:acc_1e4}) exceeds the target by roughly one order of magnitude,
which is consistent with the conservative error bound in
\cref{eq:global_error_bound}: the dense spatial revisits produce a maximum
cut depth of 75, so the accumulated truncation error from many individually
small omissions can exceed the single-edge calibration threshold.
Under the tighter $10^{-7}$ configuration (\cref{tab:acc_1e7}), all paths
remain well within the target, confirming that a single straight-path calibration transfers well to complex scan paths.

\begin{table}[h]
\centering
\caption{Observed accuracy for the configuration calibrated to a target
error of $10^{-4}$, using segment length $L_{\mathrm{seg}}=0.9$~mm and temperature
threshold $\varepsilon=5$~K.}
\label{tab:acc_1e4}
\begin{tabular}{lcc}
\hline
Path & Global max cut depth & Max relative $L^2$ error \\
\hline
Straight line        & 1  & $9.60\times10^{-5}$ \\
Bull path            & 23 & $3.08\times10^{-4}$ \\
Hybrid spiral--raster & 75 & $2.74\times10^{-3}$ \\
\hline
\end{tabular}
\end{table}

\begin{table}[h]
\centering
\caption{Observed accuracy for the configuration calibrated to a target
error of $10^{-7}$, using segment length $L_{\mathrm{seg}}=1.25$~mm and temperature
threshold $\varepsilon=0.01$~K.}
\label{tab:acc_1e7}
\begin{tabular}{lcc}
\hline
Path & Global max cut depth & Max relative $L^2$ error \\
\hline
Straight line        & 1   & $9.81\times10^{-8}$ \\
Bull path            & 81  & $8.75\times10^{-7}$ \\
Hybrid spiral--raster & 100 & $2.32\times10^{-8}$ \\
\hline
\end{tabular}
\end{table}


\subsection{Strong-Scaling Study}\label{seq:strong}

We evaluate strong scaling by fixing each scan path and increasing the
number of ranks. For each of the four scan paths in
\cref{fig:scan_paths}, we compare two partitioning strategies: uniform
partitioning, and optimized partitioning, which places partition boundaries
according to the algorithm described in \cref{seq:partitioning}. Speedup is
computed relative to the 1-rank baseline for each strategy.

We conduct two sets of experiments. In the first set, we use a grid spacing
of $h=18~\si{\micro\meter}$ with one rank per GPU, scaling from 1 to 8
ranks. In the second set, we use $h=40~\si{\micro\meter}$ and enable NVIDIA
Multi-Process Service (MPS), allowing two ranks to share a single GPU and
scaling from 1 to 16 ranks across 8 GPUs. The coarser grid produces a
smaller per-step workload that typically occupies only 30--40\% of a single
GPU, so co-locating two ranks via MPS improves device utilization.

\paragraph{One rank per GPU ($h=18~\si{\micro\meter}$)}
\cref{fig:strong_scaling} shows the results for 1--8 ranks.
At 8 ranks, the optimized partitioning achieves speedups of
$6.97\times$ (Bull), $7.35\times$ (Texas), $7.03\times$ (spiral--raster),
and $7.33\times$ (Hilbert), corresponding to 87--92\% parallel efficiency.
Uniform partitioning reaches $6.84\times$, $7.21\times$, $6.78\times$, and
$7.35\times$, respectively. The gap is most pronounced for the
spiral--raster path ($7.03\times$ vs.\ $6.78\times$), where dense spatial
revisits create regions of high cut depth that the optimized strategy avoids.


\begin{figure}[h]
    \centering
    \includegraphics[width=\linewidth]{strong_scaling_8ranks.png}
    \caption{Strong-scaling speedup for the four scan paths using
    $h=18~\si{\micro\meter}$ with one rank per GPU, scaling from 1 to 8
    ranks. Each plot compares uniform and optimized partitioning; speedup
    is measured relative to the 1-rank baseline.}
    \label{fig:strong_scaling}
\end{figure}

\paragraph{MPS ($h=40~\si{\micro\meter}$)}
\cref{fig:mps_strong_scaling} shows the results for 1--16 ranks on 8 GPUs.
At 16 ranks, optimized partitioning reaches $9.78\times$ (Bull),
$13.67\times$ (Texas), $9.70\times$ (spiral--raster), and $12.74\times$
(Hilbert). The Texas and Hilbert paths, which have longer scan paths and
therefore more segments per rank, benefit most from the additional MPS ranks.
The Bull and spiral--raster paths show diminishing returns beyond 10--12
ranks, indicating that the per-rank workload approaches the granularity at
which MPS scheduling overhead becomes visible.
Across both experiments, optimized partitioning consistently matches or
outperforms uniform partitioning, with the advantage increasing for paths
that contain regions of dense thermal dependency.


\begin{figure}[h]
    \centering
    \includegraphics[width=\linewidth]{mps_16ranks.png}
    \caption{Strong-scaling speedup for the four scan paths using
    $h=40~\si{\micro\meter}$ with MPS enabled, scaling from 1 to 16 ranks
    across 8 GPUs.}
    \label{fig:mps_strong_scaling}
\end{figure}


\subsection{Weak-Scaling Study}\label{seq:weak}

We evaluate weak scaling by increasing the number of ranks together with the
problem size so that the workload per rank remains constant. All weak-scaling
runs use a grid spacing of $h=18~\si{\micro\meter}$ with one rank per GPU,
scaling from 1 to 8 ranks.

We consider two scan-path families. The first is the hybrid spiral--raster
path, where the problem size is increased by repeating the same motif as the
rank count grows; a run with $P$ ranks uses a path with $P$ repeated motifs.
The second is the Bull path, where the horizontal extent is scaled by
$\sqrt{P}$ so that the total scanned area grows linearly with the rank
count.

\cref{fig:weak_scaling} reports the weak-scaling efficiency for both paths.
The Bull path maintains an efficiency of 1.00 across all ranks, with total
runtime remaining nearly constant at approximately 307~s. The
spiral--raster path exhibits a slight decrease from 1.00 to 0.96, with
runtime increasing from 192.2~s to 199.4~s. This is attributable to the
higher cut depths caused by the spiral's dense spatial revisits, which
increase the per-boundary correction cost relative to the simpler Bull
raster path. Overall, both paths demonstrate near-ideal weak scaling up to
8 ranks.


\begin{figure}[h]
    \centering
    \includegraphics[width=\linewidth]{figs/weak_scaling.png}
    \caption{Weak-scaling efficiency for the Bull and hybrid spiral--raster
    paths using $h=18~\mu\text{m}$ with one rank per GPU, scaling
    from 1 to 8 ranks.}
    \label{fig:weak_scaling}
\end{figure}


\subsection{Parametric Study}
\label{seq:parametric}

We investigate the effect of accuracy level on parallel performance using
the hybrid spiral--raster path. Two configurations are tested: a
low-accuracy setting with temperature threshold $\varepsilon = 5$~K and CG
tolerance $10^{-5}$, and a high-accuracy setting with $\varepsilon = 0.01$~K
and CG tolerance $10^{-10}$. The tighter CG tolerance increases the
per-solve cost accordingly.

As shown in \cref{fig:parametric_study}, under the low-accuracy
configuration, uniform and optimized partitioning achieve comparable speedup
at 8 ranks ($6.89\times$ and $7.05\times$, respectively). The larger
temperature threshold produces smaller cut depths, so the correction cost is
low regardless of where partition boundaries are placed. Under the
high-accuracy configuration, the denser thermal dependencies induced by the
tighter threshold increase the sensitivity to partition placement: optimized
partitioning achieves $7.03\times$ versus $6.78\times$ for uniform
partitioning at 8 ranks, confirming that the optimized strategy better
balances the correction workload when cross-boundary dependencies are more
numerous.

\begin{figure}[h]
    \centering
    \includegraphics[width=\linewidth]{figs/parametric_study.png}
    \caption{Strong-scaling speedup for the hybrid spiral--raster path under
    low-accuracy ($\varepsilon = 5$~K, CG tolerance $10^{-5}$) and high-accuracy ($\varepsilon = 0.01$~K, CG tolerance
    $10^{-10}$) configurations, using $h =
    18~\si{\micro\meter}$ with one rank per GPU.}
    \label{fig:parametric_study}
\end{figure}

\section{Related Work}\label{seq:related}
The computational demands of LPBF simulation have led to a clear separation between advanced spatial solvers and newer strategies for parallelizing in time. Although spatial domain decompositions effectively capture fine-scale, near-laser phenomena, they are inherently unable to bypass the sequential nature of time stepping, for which the wall-clock cost grows linearly with the total simulated time span \cite{aydin2026hermes, leonor2024go, li2023efficient}.

\subsection{Foundations of Parallel-in-Time (PinT) Integration}
Parallelization in the temporal dimension seeks to push past strong-scaling limits by revealing concurrency along the time axis. Its roots can be traced back to Nievergelt’s pioneering work \cite{nievergelt1964parallel} and the influential Parareal algorithm \cite{lions2001resolution}. Cast as a predictor–corrector scheme, Parareal partitions the global temporal evolution into multiple, independently solvable sub-interval problems \cite{lions2001resolution}. Later developments have emphasized analyzing stability in the presence of complex spectra \cite{staff2005stability} and enhancing performance via optimized scheduling strategies and manager–worker frameworks that overlap sequential and parallel stages of the computation \cite{aubanel2011scheduling}.

Iterative variants such as multigrid reduction in time  \cite{falgout2017multigrid} and PFASST \cite{emmett2012toward}, which utilizes high-order spectral deferred correction sweeps \cite{dutt2000spectral}, have further enhanced large-scale concurrency. However, these traditional PinT frameworks often struggle with LPBF due to a coarse-solver bottleneck, where it is difficult to construct a cheap model accurate enough to capture extreme thermal gradients. Furthermore, they cannot exploit the rapid spatiotemporal decay inherent to thermal diffusion.

\subsection{Temporal Strategies and Multiscale Methods for LPBF} 
Recent research has addressed the disparate scales of LPBF through specialized temporal and multiscale discretizations. Hodge \cite{hodge2021towards} proposed a multirate framework to partition fast and slow dynamics, while Viguerie et al. \cite{viguerie2022spatiotemporal} and Kopp et al. \cite{kopp2022efficient} developed spatiotemporal two-level and multi-level $hp$-finite element methods to resolve extreme spatial gradients. While effective, these methods often require intrusive modifications to the underlying numerical source code or rely on uniform temporal discretizations that do not fully exploit the path-dependent nature of the physical process.

Recently \cite{coleman2026time} achieved significant speedup using an iterative superposition of "forced" and "unforced" solutions similar to our approach. 
However, like traditional PinT methods, their methodology relies on fixed, uniform time intervals. 
Such path-blind discretization is also inherently disadvantageous for realistic scan paths, such as spirals or dense hatches, where thermal interactions are irregular.

\subsection{Novelty and Advantages of the Proposed DAG-Based Approach}\label{seq:novelty}

Our approach mainly distinguishes itself from existing iterative frameworks by replacing uniform time intervals with a path-aware, geometry-driven sparse DAG built with respect to a target accuracy. 

\paragraph{Path-Awareness and Inherent Accuracy}
Unlike standard iterative PinT methods, where accuracy depends on the number of iterations, our approach guaranties a target accuracy by construction. Using a fast analytical diffusion decay approximation (\cref{eq:point_source_decay}), we build a laser path-driven and accuracy driven dependency DAG. This allows arbitrarily complex scan paths, such as Hilbert curves or tight spirals, to be parsed into a sparse structure that explicitly exposes the underlying physical relationships. Thus, we exploit the rapid thermal decay of the LPBF process to ensure that computational resources are exclusively spent on physically significant history corrections.

\paragraph{Optimal Partitioning and Near-Ideal Scaling}
We implemented a partitioning algorithm that explicitly exploits the DAG structure. By grouping tightly coupled segments into contiguous super-segments under a budgeted cut-depth, we encapsulate dense topological dependencies internally. This minimizes bottleneck loads, resulting in the near-ideal speedup observed in our strong-scaling tests across various complex geometries on the GH200-based Vista cluster.

\paragraph{Generality and Nonlinear Extensibility}
While we demonstrate the linear case in this work, our DAG dependency logic seamlessly extends to nonlinear LPBF models. Applying this to nonlinear equations would require additional Picard-like deferred correction iterations to resolve the coupling. Crucially, these iterations would still operate over our sparse, path-aware DAG, and still avoiding the fundamental pitfalls of uniform time discretization. This non-intrusive method can be applied to any existing moving-domain solvers \cite{aydin2026hermes, leonor2024go, li2023efficient}.




\section{Conclusions}\label{seq:conclusions}
\method achieves unprecedented speedups, opening the door to algorithms for uncertainty quantification, path optimization, and integration with microstructure evolution and thermomechanical stress simulations.  
It is compatible with any moving-domain solver and involves only a small number of tunable parameters.  
Nevertheless, in its current form it has several limitations, which we outline below.


{\bf Limitations}
\begin{itemize*}

\item \emph{Constant laser parameters:}

For a laser with time-dependent parameters,  the question arises on how to define edges appropriately. 
If the laser parameters always remain constant within a single segment, the situation is simple:
the only difference from the current case is that the edge definition now depends on the specific $s_i$, $s_j$ edges. 
However, in corners and sharp turns, the laser parameters vary within a single segment.
This makes the determination of far- and near-field boundaries more involved, but we expect it to be carried out without major changes to the existing framework.

\item \emph{Linear heat transfer:}
As seen in \cite{leonor2024go, aydin2026hermes} and related references, the pronounced nonlinear behavior is restricted to the melt pool region. 
Within a multiresolution solver, the coarser levels remain linear and can thus be parallelized in time using \method. 
Moreover, localized nonlinear corrections can be applied in parallel. 
The associated algorithmic aspects are intricate and will be discussed in a separate publication.

\item \emph{Single laser:}
This is an important usage scenario. 
We consider two variants of this problem: 
Lasers operating simultaneously on different parts (a very common configuration);
and lasers operating simultaneously on the same part (less common at present, but increasingly relevant). 
\method is directly applicable to the first case because the parts are usually well separated. 
The second case is more involved, since the interaction across segments results from the superposition of multiple moving point sources. 
Extending \method to handle this important setting is, again, beyond the scope of the present paper.

\item \emph{Non-flat geometries:}
In LPBF, the process takes place on a uniform, planar powder bed, which eases the setup of numerical simulations.  
Direct energy deposition and similar methods pose greater challenges in this respect as the part surface can be curved and quite complex.
However, these techniques also involve moving point heat sources that follow line-like paths.
With appropriately chosen underlying solvers, we anticipate that \method can be extended to encompass a wider spectrum of 3D printing processes.

\end{itemize*}

% \subsection{Some Common Mistakes}\label{SCM}
% \begin{itemize}
% \item The word ``data'' is plural, not singular.
% \item The subscript for the permeability of vacuum $\mu_{0}$, and other common scientific constants, is zero with subscript formatting, not a lowercase letter ``o''.
% \item In American English, commas, semicolons, periods, question and exclamation marks are located within quotation marks only when a complete thought or name is cited, such as a title or full quotation. When quotation marks are used, instead of a bold or italic typeface, to highlight a word or phrase, punctuation should appear outside of the quotation marks. A parenthetical phrase or statement at the end of a sentence is punctuated outside of the closing parenthesis (like this). (A parenthetical sentence is punctuated within the parentheses.)
% \item A graph within a graph is an ``inset'', not an ``insert''. The word alternatively is preferred to the word ``alternately'' (unless you really mean something that alternates).
% \item Do not use the word ``essentially'' to mean ``approximately'' or ``effectively''.
% \item In your paper title, if the words ``that uses'' can accurately replace the word ``using'', capitalize the ``u''; if not, keep using lower-cased.
% \item Be aware of the different meanings of the homophones ``affect'' and ``effect'', ``complement'' and ``compliment'', ``discreet'' and ``discrete'', ``principal'' and ``principle''.
% \item Do not confuse ``imply'' and ``infer''.
% \item The prefix ``non'' is not a word; it should be joined to the word it modifies, usually without a hyphen.
% \item There is no period after the ``et'' in the Latin abbreviation ``et al.''.
% \item The abbreviation ``i.e.'' means ``that is'', and the abbreviation ``e.g.'' means ``for example''.
% \end{itemize}
% An excellent style manual for science writers is \cite{b7}.



% \subsection{Figures and Tables}
% \paragraph{Positioning Figures and Tables} Place figures and tables at the top and 
% bottom of columns. Avoid placing them in the middle of columns. Large 
% figures and tables may span across both columns. Figure captions should be 
% below the figures; table heads should appear above the tables. Insert 
% figures and tables after they are cited in the text. Use the abbreviation 
% ``Fig.~\ref{fig}'', even at the beginning of a sentence.

% \begin{table}[htbp]
% \caption{Table Type Styles}
% \begin{center}
% \begin{tabular}{|c|c|c|c|}
% \hline
% \textbf{Table}&\multicolumn{3}{|c|}{\textbf{Table Column Head}} \\
% \cline{2-4} 
% \textbf{Head} & \textbf{\textit{Table column subhead}}& \textbf{\textit{Subhead}}& \textbf{\textit{Subhead}} \\
% \hline
% copy& More table copy$^{\mathrm{a}}$& &  \\
% \hline
% \multicolumn{4}{l}{$^{\mathrm{a}}$Sample of a Table footnote.}
% \end{tabular}
% \label{tab1}
% \end{center}
% \end{table}

% \begin{figure}[htbp]
% \centerline{\includegraphics{fig1.png}}
% \caption{Example of a figure caption.}
% \label{fig}
% \end{figure}

% Figure Labels: Use 8 point Times New Roman for Figure labels. Use words 
% rather than symbols or abbreviations when writing Figure axis labels to 
% avoid confusing the reader. As an example, write the quantity 
% ``Magnetization'', or ``Magnetization, M'', not just ``M''. If including 
% units in the label, present them within parentheses. Do not label axes only 
% with units. In the example, write ``Magnetization (A/m)'' or ``Magnetization 
% \{A[m(1)]\}'', not just ``A/m''. Do not label axes with a ratio of 
% quantities and units. For example, write ``Temperature (K)'', not 
% ``Temperature/K''.

% %\section*{Acknowledgment}

% \section*{References}

% Please number citations consecutively within brackets \cite{b1}. The 
% sentence punctuation follows the bracket \cite{b2}. Refer simply to the reference 
% number, as in \cite{b3}---do not use ``Ref. \cite{b3}'' or ``reference \cite{b3}'' except at 
% the beginning of a sentence: ``Reference \cite{b3} was the first $\ldots$''

% Number footnotes separately in superscripts. Place the actual footnote at 
% the bottom of the column in which it was cited. Do not put footnotes in the 
% abstract or reference list. Use letters for table footnotes.

% Unless there are six authors or more give all authors' names; do not use 
% ``et al.''. Papers that have not been published, even if they have been 
% submitted for publication, should be cited as ``unpublished'' \cite{b4}. Papers 
% that have been accepted for publication should be cited as ``in press'' \cite{b5}. 
% Capitalize only the first word in a paper title, except for proper nouns and 
% element symbols.

% For papers published in translation journals, please give the English 
% citation first, followed by the original foreign-language citation \cite{b6}.

% \begin{thebibliography}{00}
% \bibitem{b1} G. Eason, B. Noble, and I. N. Sneddon, ``On certain integrals of Lipschitz-Hankel type involving products of Bessel functions,'' Phil. Trans. Roy. Soc. London, vol. A247, pp. 529--551, April 1955.
% \bibitem{b2} J. Clerk Maxwell, A Treatise on Electricity and Magnetism, 3rd ed., vol. 2. Oxford: Clarendon, 1892, pp.68--73.
% \bibitem{b3} I. S. Jacobs and C. P. Bean, ``Fine particles, thin films and exchange anisotropy,'' in Magnetism, vol. III, G. T. Rado and H. Suhl, Eds. New York: Academic, 1963, pp. 271--350.
% \bibitem{b4} K. Elissa, ``Title of paper if known,'' unpublished.
% \bibitem{b5} R. Nicole, ``Title of paper with only first word capitalized,'' J. Name Stand. Abbrev., in press.
% \bibitem{b6} Y. Yorozu, M. Hirano, K. Oka, and Y. Tagawa, ``Electron spectroscopy studies on magneto-optical media and plastic substrate interface,'' IEEE Transl. J. Magn. Japan, vol. 2, pp. 740--741, August 1987 [Digests 9th Annual Conf. Magnetics Japan, p. 301, 1982].
% \bibitem{b7} M. Young, The Technical Writer's Handbook. Mill Valley, CA: University Science, 1989.
% \end{thebibliography}
% \vspace{12pt}
% \color{red}
% IEEE conference templates contain guidance text for composing and formatting conference papers. Please ensure that all template text is removed from your conference paper prior to submission to the conference. Failure to remove the template text from your paper may result in your paper not being published.

% \begin{thebibliography}{00}
% \bibitem{b1} J.~Nievergelt, ``Parallel methods for integrating ordinary differential equations,'' \emph{Communications of the ACM}, vol. 7, no. 12, pp. 731--733, 1964.
% \bibitem{b2} J.-L. Lions, Y. Maday, and G. Turinici, ``A `parareal' in time discretization of PDE's,'' \emph{C. R. Acad. Sci. Paris, Ser. I}, vol. 332, no. 7, pp. 661--668, 2001.
% \end{thebibliography}
% \vspace{12pt}


\bibliographystyle{IEEEtran}
\bibliography{refs}
\end{document}
